\chapter{Азбука Морзе}
% полная версия: chapters/04_code.tex

В азбуке Морзе два знака --- точка и тире --- и паузы между ними. У нейрона алфавит ещё беднее: знак один, короткий щелчок. Значит, всё сообщение записано в паузах. На каком же языке разговаривают нейроны?

\section*{Сколько можно сказать щелчками}

Начнём как инженер связи: сколько информации вообще помещается в таком канале? Если получатель различает время прихода щелчка с точностью до миллисекунды, один щелчок может нести около восьми бит. Если же получатель только считает, сколько щелчков пришло за секунду, --- из того же потока он извлечёт в десятки раз меньше. Почти вся ёмкость канала сидит во \emph{времени}.

\pencilfig{4}{\pencilart{4}}{Сверху азбука Морзе, снизу нейрон: знак один, щелчок; всё сообщение записано в паузах.}

Но канал шумный. Сам нейрон, как ни странно, точен: если много раз подать на него один и тот же меняющийся ток, он будет щёлкать в одни и те же моменты с точностью до миллисекунды. Шумят соединения: больше половины щелчков, добежавших до конца провода, просто теряется --- вещество выбрасывается не каждый раз. В живой коре поток щелчков выглядит почти случайным. Шум ли это --- или сообщение, которое мы пока не умеем прочитать? Хорошо сжатое сообщение для постороннего неотличимо от случайного, так что вопрос открыт.

\section*{Медленно, но надёжно}

Самый простой код --- считать щелчки: чем чаще, тем больше число. Этому коду не страшны ни потери, ни дрожание. Но он медленный. Чтобы узнать частоту с точностью в десять процентов, нужна сотня щелчков, а это секунды. Между тем вы узнаёте животное на картинке, показанной на мгновение, примерно за полторы десятых секунды, и сигнал при этом проходит около десятка ступеней обработки. На каждой ступени каждый нейрон успевает щёлкнуть от силы раз.

Выход --- взять много нейронов сразу: не один слушатель считает щелчки минуту, а сотня слушателей --- мгновение. Только и тут подвох: шумы соседних нейронов немного похожи, как слухи в толпе. Усреднение по нескольким десяткам соседей помогает, а дальше почти перестаёт.

\pencilfig{122}{%
% error of the group-average rate relative to one neuron, sqrt(c + (1-c)/N): c = 0.1 (neighbours' noise
% is slightly alike; full edition, chapter 4) and c = 0 (independent noise). x = 0.8 log2 N, y = 2.2 * error
\draw[pen] (0,0) -- (5.9,0); \draw[pen] (0,0) -- (0,2.45);
\foreach \x/\n in {0/1, 2.66/10, 5.31/100}{\draw[pcGray, line width=0.5pt] (\x,0) -- (\x,-0.08); \node[lbl, anchor=north] at (\x,-0.08) {\n};}
\node[lbl, anchor=north] at (2.95,-0.38) {сколько нейронов усредняем};
\node[lbl, anchor=south, rotate=90] at (-0.1,1.2) {ошибка};
\draw[pcGray, densely dashed, line width=0.5pt] (0,0.70) -- (5.6,0.70);
\draw[penlite=pcBlue] plot[smooth] coordinates {(0.00,2.20) (0.47,1.80) (0.80,1.56) (1.27,1.27) (1.60,1.10) (2.07,0.90) (2.40,0.78) (2.87,0.64) (3.20,0.55) (3.67,0.45) (4.00,0.39) (4.47,0.32) (4.80,0.28) (5.27,0.22) (5.60,0.19)};
\draw[penlite=pcRed] plot[smooth] coordinates {(0.00,2.20) (0.47,1.84) (0.80,1.63) (1.27,1.39) (1.60,1.25) (2.07,1.10) (2.40,1.01) (2.87,0.92) (3.20,0.87) (3.67,0.82) (4.00,0.79) (4.47,0.76) (4.80,0.74) (5.27,0.73) (5.60,0.72)};
\node[lbl, text=pcRed, anchor=west, align=left] at (5.7,0.74) {похожие шумы,\\как в коре};
\node[lbl, text=pcBlue, anchor=west] at (5.7,0.19) {независимые шумы};
\node[lbl, anchor=west] at (0.9,2.15) {один нейрон};
\node[lbl, anchor=north west] at (0.05,0.66) {втрое меньше};
}{Ошибка частоты, усреднённой по группе нейронов. Будь шумы независимы, она падала бы и дальше. Похожие шумы соседей останавливают её примерно на трети, и этот предел достигнут уже на нескольких десятках нейронов.}

\section*{Быстро: кто первый}

Второй выход --- писать число во времени. Сильный сигнал заставит нейрон щёлкнуть рано, слабый --- поздно. Один-единственный щелчок уже несёт число. Но от чего отсчитывать время, если нет часов? Можно сравнивать нейроны друг с другом: кто щёлкнул первым, кто вторым. Порядок прихода щелчков десяти нейронов несёт больше двадцати бит. А ещё отметкой начала служит общий залп --- как длинная пауза между словами в азбуке Морзе.

\pencilfig{121}{%
% latency code: one click per neuron; the stronger the input, the earlier the click
\foreach \y/\t/\l/\k in {1.5/1.1/{сильный вход}/1, 0.95/2.4/{средний}/2, 0.4/4.2/{слабый}/3}{
  \draw[pen] (0.4,\y) -- (5.1,\y);
  \draw[pcBlue, line width=0.9pt] (\t,\y) -- (\t,\y+0.42);
  \node[lbl, anchor=east] at (0.3,\y+0.1) {\l};
  \draw[dotpen=pcAmber, wash=pcAmber] (5.6,\y+0.16) circle (0.17);
  \node[font=\scriptsize, text=pcAmber!60!black] at (5.6,\y+0.16) {\k};}
\draw[pcGray, densely dashed, line width=0.5pt] (0.55,0.25) -- (0.55,2.1);
\node[lbl, anchor=south] at (0.55,2.1) {раздражитель};
\node[lbl, text=pcAmber!70!black, anchor=south] at (5.6,2.1) {порядок};
\draw[arr] (0.7,0.05) -- (4.9,0.05); \node[lbl, anchor=north] at (2.8,0.02) {время};
}{Сильный вход --- ранний щелчок, слабый --- поздний. Порядок щелчков говорит, какой вход сильнее, даже если момент начала неизвестен.}

\section*{Код выбирает слушатель}

Одна и та же последовательность щелчков несёт и частоту, и времена, и порядок. Что из этого будет прочитано, решает получатель. Нейрон с «широкой дыркой» в ведре медленно копит воду и слышит только частоту. Нейрон с быстрым забыванием слышит только совпадения. А тормоза, как мы видели, умеют переключать нейрон из одного режима в другой.

Даже соединения фильтруют сигнал. Одни быстро устают и потому сообщают не частоту, а её \emph{изменение}. Другие, наоборот, пропускают лучше пачки из нескольких щелчков подряд: у нейрона появляется «тире». Тормозные нейроны тем временем расставляют знаки препинания --- режут поток на окна и приглушают его, когда он слишком громок.

\section*{Экономный язык}

И последнее: этот язык дорог. Каждый щелчок стоит энергии, а весь мозг потребляет около двадцати ватт. Если разделить эти ватты на все нейроны, получится в среднем порядка одного щелчка в секунду на нейрон. Большинство нейронов большую часть времени молчат. Код мозга обязан быть скупым.

В азбуке Морзе самая частая буква английского языка, E, --- это одна точка: частое кодируют коротко. У нейронов похоже: то, что ожидаемо, стоит мало щелчков. Постоянный раздражитель постепенно перестаёт вызывать ответ, датчики сообщают о переменах, а дофамин --- только о сюрпризах. Мозг тратит свои щелчки на новости.

\fullversion{4}
